Amorphous and network materials exhibit complex yielding and fracture behaviour under deformation, arising from the interplay between disorder, stress redistribution, and local failure processes. Understanding how these microscopic mechanisms give rise to macroscopic yielding and fracture remains a key challenge. In this thesis, we study yielding and fracture in simplified models of disordered soft materials across three separate problems.

First, we consider an elastoplastic model of an amorphous material under cyclic deformation. Depending on the imposed strain amplitude and degree of annealing, the system evolves either towards an absorbing state, in which plastic activity ceases, or towards yielding and failure through sustained irreversible rearrangements. Near this transition, the dynamics become long-lived and intermittent. For sufficiently low annealing (initial disorder), we show that activity within the system becomes dilute, with local plastic yielding events distributed across the system in a dilute way, allowing the late-time dynamics to be interpreted in terms of weakly interacting elements and the approach to absorption as a stochastic first-passage process.

Second, we study fracture in disordered fibre networks subjected to a step strain using a mesoscale network model with thermally activated bond breaking. Rather than relaxing smoothly or failing immediately, the networks can remain apparently stable for long times before displaying an abrupt, system-spanning fracture. This ultra-delayed failure is driven by bonds strained just below their rupture threshold, which break via thermal activation and trigger cascade-like avalanches that culminate in macroscopic failure.

Finally, we investigate the mechanical response of double-network materials composed of a stiff, brittle sacrificial network coupled to a softer, ductile matrix. Using a mesoscale network model that resolves both macroscopic fracture and local bond failure, we show that load sharing between the two networks delocalises stress and suppresses stress concentration around damaged regions. This inhibits cascade-like bond failure and prevents crack nucleation, such that brittle fracture is replaced by diffuse microcracking. As a result, the material combines the stiffness of the sacrificial network with the ductility of the matrix, leading to high strength and large failure strains.